% ======================================================================
\section{Unitarity relations: derivation}\label{app:unitarity}
% ======================================================================

In this appendix the relations among the coefficients of $\Omega$ that follow from $\Omega^\dagger\Omega=1$ are derived order by order in $g$.

The conjugate of Eq.~\eqref{eq:OmegaExp} is
\begin{align}\label{eq:OmegaDag}
\Omega^\dagger = \ &\mf{N}^\dagger
+\int\frac{d^3p}{(2\pi)^3}\Big[\mf{A}^{\dagger b}_{1j}(p)\,a^{\dagger b}_j(p)+\mf{A}^{\dagger b}_{2j}(p)\,a^{b}_j(p)\Big]
\notag\\&
+\int\frac{d^3p}{(2\pi)^3}\frac{d^3q}{(2\pi)^3}\Big[\mf{B}^{\dagger cb}_{1kj}(p,q)\,a^{\dagger b}_j(p)a^{\dagger c}_k(q)+\mf{B}^{\dagger cb}_{2kj}(p,q)\,a^{b}_j(p)a^{c}_k(q)
\notag\\&\hspace{3.4cm}+\mf{B}^{\dagger cb}_{3kj}(p,q)\,a^{\dagger b}_j(p)a^{c}_k(q)\Big]
\notag\\&
+\int\frac{d^3p}{(2\pi)^3}\frac{d^3q}{(2\pi)^3}\Big[\mf{C}^{\dagger dcb}_{1lkj}(p,q)\,a^{\dagger b}_j(p)a^{c}_k(p-q)a^{d}_l(q)
\notag\\&\hspace{3.4cm}+\mf{C}^{\dagger dcb}_{2lkj}(p,q)\,a^{\dagger d}_l(q)a^{\dagger c}_k(p-q)a^{b}_j(p)\Big]+\dots,
\end{align}
where the dagger on a coefficient is the Hermitian conjugate on $\mathcal H_{\rm val}$ and the indices and arguments are those of the coefficient that is conjugated. The product $\Omega^\dagger\Omega$ is brought to normal order with the commutator~\eqref{eq:aadagger}. The coefficients commute with $a$ and $a^\dagger$, but not with each other: in each product the coefficient that comes from $\Omega^\dagger$ stands to the left of the one from $\Omega$. The normal-ordered products of $a^\dagger$ and $a$ are linearly independent, so that the coefficient of each of them has to vanish, except for the identity. We give the relations in the order in which they arise; the coefficients of the conjugate structures give the conjugate relations.

\paragraph{Order $g$.} The terms of order $g$ in $\Omega^\dagger\Omega$ are linear in the coefficients $\mf A^{(1)}_{1,2}$ and $\mf C_{1,2}$. The coefficient of $a^{\dagger b}_j(p)$ is $\mf A^{(1)b}_{2j}(p)+\mf A^{(1)\dagger b}_{1j}(p)$, and the coefficient of $a^{\dagger d}_l(q)\,a^{\dagger c}_k(p-q)\,a^b_j(p)$ is $\mf C^{dcb}_{1lkj}(p,q)+\mf C^{\dagger dcb}_{2lkj}(p,q)$. Hence
\begin{align}
\mf A^{(1)\dagger b}_{1j}(p)&=-\mf A^{(1)b}_{2j}(p)\,,\label{eq:A11unit}\\
\mf C^{\dagger dcb}_{2lkj}(p,q)&=-\mf C^{dcb}_{1lkj}(p,q)\,.\label{eq:C2unit}
\end{align}
At this order $\Omega=1+\Omega^{(1)}$ with $\Omega^{(1)}$ anti-Hermitian. Eq.~\eqref{eq:A11unit} is satisfied by the results of the diagrammatic calculation, Eqs.~\eqref{eq:A1mom} and~\eqref{eq:A11mom}. Eq.~\eqref{eq:C2unit} determines $\mf C_2$, the amplitude for the merging of two gluons, which we have not computed directly.

\paragraph{Order $g^2$.} At this order the terms linear in $\mf N_2$ and $\mf B_{1,2,3}$ are accompanied by the products of two coefficients of order $g$, normal ordered. The coefficient of the identity receives $\mf N_2+\mf N_2^\dagger$ and the contraction of $\mf A^{(1)\dagger}_2a$ with $\mf A^{(1)}_2a^\dagger$:
\begin{align}\label{eq:Nunit}
\mf N_2+\mf N_2^\dagger=-\int\frac{d^3p}{(2\pi)^3}\,\mf A^{(1)\dagger b}_{2j}(p)\,\mf A^{(1)b}_{2j}(p)\,.
\end{align}
Unitarity fixes the Hermitian part of $\mf N_2$ only. With the canonical normalisation of Section~\ref{sec:lcpt}, $\mf N_2$ is Hermitian, and
\begin{align}\label{eq:N2unit}
\mf N_2=-\frac12\int\frac{d^3p}{(2\pi)^3}\,\mf A^{(1)\dagger b}_{2j}(p)\,\mf A^{(1)b}_{2j}(p)
=-\frac{g^2}{8\pi^3}\int_\Lambda^\vee\frac{dp^+}{p^+}\int d^2p\,\frac{\rho^b(\v p)\,\rho^b(-\v p)}{\v p^2}\,,
\end{align}
where Eq.~\eqref{eq:A1mom} was inserted. This reproduces Eq.~\eqref{eq:N2mom}: the coefficient of the identity is minus one half of the probability to find one soft gluon. The coefficient of $a^{\dagger b}_j(p)\,a^{\dagger c}_k(q)$, symmetrised according to Eq.~\eqref{eq:symm}, gives
\begin{align}\label{eq:B1unit}
\mf B^{\dagger cb}_{1kj}(p,q)=\ &-\mf B^{cb}_{2kj}(p,q)+\mf C^{cbd}_{1kjl}(p+q,q)\,\mf A^{(1)d}_{2l}(p+q)\notag\\
&+\frac12\Big[\mf A^{(1)b}_{2j}(p)\,\mf A^{(1)c}_{2k}(q)+\mf A^{(1)c}_{2k}(q)\,\mf A^{(1)b}_{2j}(p)\Big]\,,
\end{align}
where Eqs.~\eqref{eq:A11unit} and~\eqref{eq:C2unit} were used. It determines $\mf B_1$, the amplitude for the absorption of two gluons. The second term is the absorption of one gluon by the valence charge after the two gluons have merged; the last term is the symmetrised product of two absorptions. Finally, the coefficient of $a^{\dagger b}_j(p)\,a^c_k(q)$ gives
\begin{align}\label{eq:B3unit}
&\mf B^{\dagger cb}_{3kj}(p,q)+\mf B^{bc}_{3jk}(q,p)+\mf A^{(1)b}_{2j}(p)\,\mf A^{(1)\dagger c}_{2k}(q)+\mf A^{(1)\dagger c}_{2k}(q)\,\mf A^{(1)b}_{2j}(p)\notag\\
&\quad+2\,\mf A^{(1)\dagger d}_{2l}(q-p)\,\mf C^{bdc}_{1jlk}(q,p)+2\,\mf C^{\dagger cdb}_{1klj}(p,q)\,\mf A^{(1)d}_{2l}(p-q)\notag\\
&\quad+2\,(2\pi)^3\delta^{(3)}(p-q)\int\frac{d^3r}{(2\pi)^3}\,\mf C^{\dagger edb}_{1mlj}(p,r)\,\mf C^{edc}_{1mlk}(p,r)=0\,.
\end{align}
The first term of the second line is present for $q^+>p^+$ and the second for $p^+>q^+$. Eq.~\eqref{eq:B3unit} contains only coefficients that are fixed by the diagrammatic calculation. It is a check of $\mf B_3$, and it is satisfied by Eq.~\eqref{eq:B3mom} term by term:
\begin{itemize}
\item The instantaneous term and the commutator term of Eq.~\eqref{eq:B3mom} cancel between $\mf B_3$ and $\mf B_3^\dagger$, as shown in Section~\ref{sec:summarywf}.
\item The anticommutator term gives
\begin{align}
\mf B^{bc}_{3jk}(q,p)+\mf B^{\dagger cb}_{3kj}(p,q)=-\frac{2g^2\,p^jq^k}{\v p^2\v q^2\sqrt{p^+q^+}}\,\big\{\rho^b(-\v p),\rho^c(\v q)\big\}\,,
\end{align}
which is minus the sum of the two products of $\mf A^{(1)}_2$ in the first line, by Eq.~\eqref{eq:A1mom}.
\item The terms with the three-gluon vertex cancel the second line. This fixes, in particular, the relative sign of the term $\propto\Theta_>$ of Eq.~\eqref{eq:B3mom} and $\mf C_1$.
\item The last line is cancelled by the $\rho$-independent term~\eqref{eq:B3rho0}.
\end{itemize}
The origin of these cancellations is an identity for the second-order terms of Eq.~\eqref{eq:pt2}. For two states $i\neq f$ with $E_i\neq E_f$ and an intermediate state $m$,
\begin{align}\label{eq:adaggera_identity}
\frac{V_{fm}V_{mi}}{(E_i-E_f)(E_i-E_m)}+\Big[\frac{V_{im}V_{mf}}{(E_f-E_i)(E_f-E_m)}\Big]^\dagger=-\Big[\frac{V_{mf}}{E_f-E_m}\Big]^\dagger\,\frac{V_{mi}}{E_i-E_m}\,,
\end{align}
which follows from $V_{mf}^\dagger=V_{fm}$ and $\frac{1}{E_i-E_m}-\frac{1}{E_f-E_m}=\frac{E_f-E_i}{(E_i-E_m)(E_f-E_m)}$, with the order of the operators kept. The left-hand side is the contribution of the state $m$ to $\mf B_3+\mf B_3^\dagger$, the right-hand side is minus the product of two first-order amplitudes. With the vacuum as intermediate state the right-hand side is the product of two absorption amplitudes, with a two-gluon intermediate state it is the product of two emission amplitudes or of an emission and a splitting amplitude. The pole at $E_i=E_f$ cancels on the left-hand side. Eq.~\eqref{eq:adaggera_identity} holds point by point for $E_i\neq E_f$ and therefore for the principal-value prescription~\eqref{eq:PV}.

\paragraph{Order $g^3$.} At this order we need only the coefficient of $a^{\dagger b}_j(p)$. It contains $\mf A^{(3)}_2+\mf A^{(3)\dagger}_1$, the products of $\mf N_2$ with the coefficients of order $g$, and the contractions of $\mf A^{(1)\dagger}_2a$ with $\mf B_2a^\dagger a^\dagger$, of $\mf B_3^\dagger a^\dagger a$ with $\mf A^{(1)}_2a^\dagger$, and of $\mf C_1^\dagger a^\dagger aa$ with $\mf B_2a^\dagger a^\dagger$; the result is Eq.~\eqref{eq:A31unit} of the main text.
The commutator arises as follows: the two terms with $\mf N_2$ are $\mf N_2^\dagger\,\mf A^{(1)}_2$, from $\mf N^\dagger$ in $\Omega^\dagger$ and $\mf A_2a^\dagger$ in $\Omega$, and $\mf A^{(1)\dagger}_1\,\mf N_2=-\mf A^{(1)}_2\,\mf N_2$, from $\mf A_1^\dagger a^\dagger$ in $\Omega^\dagger$ and $\mf N$ in $\Omega$. Eq.~\eqref{eq:A31unit} determines $\mf A^{(3)}_1$, the amplitude for the absorption of a gluon at order $g^3$, in terms of coefficients that are fixed by the diagrammatic calculation. In the remaining relations of order $g^3$ the coefficients denoted by dots in Eq.~\eqref{eq:OmegaExp} appear; they are not needed.

\paragraph{The order of the operators.} For commuting charges the commutator in Eq.~\eqref{eq:A31unit} is absent. Here it is not: $\mf N_2$ is bilinear in the charges at two different points, and with Eqs.~\eqref{eq:rhoalgebra}, \eqref{eq:N2mom} and~\eqref{eq:A1mom}
\begin{align}\label{eq:commN2A}
\big[\mf N_2,\mf A^{(1)a}_{2i}(k)\big]=-\log\frac{\vee}{\Lambda}\ \frac{\sqrt2\,g^3\,k^i}{8\pi^3\sqrt{k^+}\,\v k^2}\int d^2p\ \frac{if^{abc}\,\{\rho^b(-\v p),\rho^c(-\v k+\v p)\}}{\v p^2}\,.
\end{align}
This is a term with two charges, of the same order in $g$ as the terms with three charges from which it originates. Comparing with Eq.~\eqref{eq:A3tworho}, it is equal to the term $-2k^i\log(\vee/\Lambda)$ of $[\mf A^{(3)}_2]_{\rho\rho}$. Two consequences follow. First, the order of the charges in the three-$\rho$ term~\eqref{eq:A3threerho} is not a matter of convention: writing it as $\mf N_2\,\mf A^{(1)}_2$ instead of $\mf A^{(1)}_2\,\mf N_2$ would change $\mf A^{(3)}_2$ by the two-$\rho$ term~\eqref{eq:commN2A}. Second, with the order of Eq.~\eqref{eq:A3threerho} the first two terms of Eq.~\eqref{eq:A31unit} combine into Eq.~\eqref{eq:N2Acomb},
so that the three-$\rho$ parts of the emission and of the absorption amplitude are $\mf A^{(1)}_2\mf N_2$ and $-(\mf N_2\,\mf A^{(1)}_2)^\dagger=\mf A^{(1)}_1\,\mf N_2$, respectively. In both, the amplitude of order $g$ stands to the left of the norm factor, in agreement with Eq.~\eqref{eq:pt3} and its analogue for the one-gluon state.


\paragraph{Verification.} The relations of this appendix were verified by a symbolic normal ordering of $\Omega^\dagger\Omega$ with non-commuting coefficients for three discrete modes (script \texttt{scripts/check\_unitarity.py}): after the substitution of Eqs.~\eqref{eq:A11unit}, \eqref{eq:C2unit}, \eqref{eq:N2unit}, \eqref{eq:B1unit} and~\eqref{eq:A31unit} the coefficients of $1$, $a^\dagger$, $a^\dagger a^\dagger$ and, at order $g$, $a^\dagger a^\dagger a$ vanish identically, and the coefficient of $a^\dagger a$ is Eq.~\eqref{eq:B3unit}. The triple-vertex part of Eq.~\eqref{eq:B3unit} was checked numerically with Eqs.~\eqref{eq:B3mom}, \eqref{eq:C1mom} and~\eqref{eq:A1mom} (relative residual $10^{-16}$; with the opposite sign of the term $\propto\Theta_>$ the residual is of order one), and Eq.~\eqref{eq:commN2A} with $SU(2)$ charges on a lattice (script \texttt{scripts/check\_misc.py}).
